\section{Method}\label{sec:method}

All mechanisms studied here reduce the search space of \eqref{eq:obj}--\eqref{eq:int} in one of
two ways. \emph{Facility restriction} keeps a subset $K\subseteq I$ and solves the problem with
$y_i=0$ for $i\notin K$ (Sections~\ref{sec:gnn}--\ref{sec:kernel}). \emph{Customer restriction}
keeps the current assignment of all customers outside a subset $F\subseteq J$ and re-optimizes
only $F$ (Sections~\ref{sec:clns}--\ref{sec:memory}). The hybrid (Section~\ref{sec:hybrid})
applies the first and then the second. For each mechanism we give a learned and a learning-free
variant, so that the contribution of learning can be isolated.

\subsection{Facility scores}\label{sec:gnn}
A \emph{score} is a vector $p\in\mathbb{R}^m$; facilities are ranked by decreasing $p_i$.

\paragraph{LP score (no learning).} With $(y^{\mathrm{LP}},r)$ the opening values and reduced
costs of the strong LP (Section~\ref{sec:lp}),
\begin{equation}\label{eq:lpscore}
p^{\mathrm{LP}}_i = y^{\mathrm{LP}}_i - \epsilon\,\frac{r_i}{\max_{k\in I}|r_k|},\qquad
\epsilon=10^{-3},
\end{equation}
so that facilities are ordered by their LP opening value and ties among facilities closed in the
LP ($y^{\mathrm{LP}}_i=0$) are broken by increasing reduced cost, i.e., by how little the LP bound
would deteriorate if they were forced open.

\paragraph{Classical score (no learning, no LP).} The service ratio of the greedy \emph{add}
heuristic: let $j_1,j_2,\dots$ be the customers sorted by increasing $u_{ij}$ and $q$ the largest
index with $\sum_{k\le q} d_{j_k}\le Q_i$; then
\begin{equation}
p^{\mathrm{cl}}_i=-\frac{f_i+\sum_{k\le q}c_{ij_k}}{\sum_{k\le q}d_{j_k}},
\end{equation}
the negative of
the average cost per unit of demand when $i$ serves its cheapest customers up to capacity.

\paragraph{GNN score (learned).} Each instance is a bipartite graph $G=(I\cup J,E)$, where $E$
contains the pair $(i,j)$ if $i$ is among the $k=10$ cheapest facilities of $j$ or $j$ is among
the $k=10$ cheapest customers of $i$ (by $u_{ij}$). Node and edge features are normalized by
statistics of the instance itself, so that one model applies across sizes: for facilities,
capacity and fixed cost relative to the instance means, the classical score, local demand
density and a profile-based competition count; for customers, demand, cheapest unit cost and
regret; for edges, the unit cost $u_{ij}$ and its rank from both endpoints.
Initial embeddings are $h^{0}_i=\phi_I(\mathbf{x}_i)$, $h^{0}_j=\phi_J(\mathbf{x}_j)$,
$e_{ij}=\phi_E(\mathbf{x}_{ij})$ in $\mathbb{R}^{64}$. Each of $R=3$ rounds updates customers
from facilities and then facilities from customers:
{\small
\begin{align}
\mu^{t}_{ij} &= \psi^{t}_{IJ}\bigl([h^{t}_i,e_{ij}]\bigr), &
m^{t}_{j} &= \Bigl[\tfrac{1}{|N(j)|}\textstyle\sum_{i\in N(j)}\mu^{t}_{ij},\
\tfrac{1}{\Delta_J}\sum_{i\in N(j)}\mu^{t}_{ij}\Bigr], \nonumber\\
&& h^{t+1}_j &= \mathrm{LN}\bigl(h^{t}_j+\gamma^{t}_{J}([h^{t}_j,m^{t}_j])\bigr), \label{eq:mpj}\\
\nu^{t}_{ji} &= \psi^{t}_{JI}\bigl([h^{t+1}_j,e_{ij}]\bigr), &
m^{t}_{i} &= \Bigl[\tfrac{1}{|N(i)|}\textstyle\sum_{j\in N(i)}\nu^{t}_{ji},\
\tfrac{1}{\Delta_I}\sum_{j\in N(i)}\nu^{t}_{ji}\Bigr], \nonumber\\
&& h^{t+1}_i &= \mathrm{LN}\bigl(h^{t}_i+\gamma^{t}_{I}([h^{t}_i,m^{t}_i])\bigr), \label{eq:mpi}
\end{align}}
where $N(\cdot)$ is the neighborhood in $G$, $\Delta_J$ and $\Delta_I$ are the maximum degrees
(so the second aggregate is a sum normalized to $[0,1]$, which keeps degree information that the
mean discards), $\psi$, $\gamma$ and $\phi$ are two-layer perceptrons, $[\cdot,\cdot]$ is
concatenation and $\mathrm{LN}$ is layer normalization. The output is one logit per facility,
$p^{\mathrm{GNN}}_i=\eta(h^{R}_i)$. Conditioning messages on $e_{ij}$ matters because the
relevant information (how expensive $i$ is for $j$ relative to $j$'s alternatives) lives on the
edges.

The target is a \emph{soft} label that tolerates alternative optima. Let $\mathcal{P}$ be the
pool of solutions found for a training instance by SCIP (60~s) and LNS (30~s), and
$\mathcal{P}_{\tau}=\{a\in\mathcal{P}: z(a)\le(1+\tau)\min_{a'\in\mathcal{P}}z(a')\}$ with
$\tau=1\%$. The label of facility $i$ is the frequency with which it is open in near-best
solutions,
\begin{equation}\label{eq:label}
q_i=\frac{1}{|\mathcal{P}_{\tau}|}\sum_{a\in\mathcal{P}_{\tau}}\mathbf{1}[\,i\in a(J)\,]\in[0,1],
\end{equation}
and the loss is the class-balanced binary cross-entropy
\begin{equation}\label{eq:loss}
\mathcal{L}=-\frac{1}{m}\sum_{i\in I}\omega_i\Bigl[q_i\log\sigma(p_i)+(1-q_i)\log\bigl(1-\sigma(p_i)\bigr)\Bigr],
\qquad
\omega_i=\begin{cases}\dfrac{1}{2\pi} & q_i>\tfrac12,\\[2mm] \dfrac{1}{2(1-\pi)} & \text{otherwise,}\end{cases}
\end{equation}
with $\sigma$ the logistic function and $\pi=\min\{0.98,\max\{0.02,\frac1m\sum_i q_i\}\}$ the
mean label of the instance; the weights $\omega_i$ give open and closed facilities the same total
weight, since typically fewer than half of the facilities are open. We train three seeds with
early stopping on validation instances.

\subsection{Restricted problems and feasibility repair}\label{sec:repair}
For $K\subseteq I$, the restricted problem $\mathrm{P}(K)$ is \eqref{eq:obj}--\eqref{eq:int} with
$y_i=0$ for $i\notin K$; its optimal value satisfies $z^\ast(K)\ge z^\ast$, with equality iff
some optimal solution opens only facilities of $K$. Given a ranking and a target fraction
$\rho\in(0,1]$, we take the $\lceil\rho m\rceil$ best facilities and \emph{repair} the set in two
steps, adding facilities in ranking order:
\begin{enumerate}
\item \emph{Packing.} Add facilities until the first-fit-decreasing heuristic (customers by
decreasing demand, each placed in the first kept facility with enough residual capacity) packs
all demands. This is a sufficient condition for $\mathrm{P}(K)$ to be feasible; the necessary
condition $\sum_{i\in K}Q_i\ge\sum_j d_j$ is not sufficient under single sourcing.
\item \emph{Coverage.} For every customer with none of its five cheapest facilities in $K$, add
the best-ranked of those five.
\end{enumerate}
We denote the result $\mathrm{rep}(\rho)$. On our instances the repair alone keeps at least 41\%
of the facilities, so a nominal $\rho=0.2$ does not yield a problem five times smaller; this
lower bound on the kept fraction is a property of capacitated instances, not of the ranking.

\subsection{Adaptive expansion with a reduced-cost certificate}\label{sec:expansion}
Let $\rho_1<\rho_2<\rho_3<\rho_4$ be $0.2,0.4,0.6,1$ and $K_s=\mathrm{rep}(\rho_s)\cup a^{s-1}(J)$,
where $a^{s-1}$ is the incumbent after stage $s-1$ (so the incumbent is always feasible for the
next stage) and $K_4=I$. Stage $s$ solves $\mathrm{P}(K_s)$ with SCIP, warm-started with
$a^{s-1}$, for a quarter of the remaining budget (the last stage receives all of it).

\begin{proposition}[Reduced-cost exclusion]\label{prop:rc}
Let $L$ be the optimal value of the strong LP and $r_i\ge0$ the reduced cost \eqref{eq:rc} of
$y_i$ at an optimal dual solution. Every feasible solution of
\eqref{eq:obj}--\eqref{eq:int} with $y_i=1$ has cost at least $L+r_i$. Consequently, if
$L+r_i\ge U$ for an incumbent of cost $U$, no solution that opens $i$ is better than the
incumbent.
\end{proposition}
\begin{proof}
Let $(v,w,s,t)$ be an optimal dual solution and $(x,y)$ any LP-feasible point. Substituting
$f_k=r_k+Q_kw_k+\sum_j s_{kj}-t_k$ and $c_{kj}=\bar c_{kj}+v_j-d_jw_k-s_{kj}$ from \eqref{eq:rc}
into the objective and regrouping by multiplier,
\begin{align*}
\sum_{k} f_k y_k+\sum_{k,j} c_{kj}x_{kj}
={}& \sum_k r_ky_k+\sum_{k,j}\bar c_{kj}x_{kj}+\sum_j v_j\underbrace{\sum_k x_{kj}}_{=1}
+\sum_k w_k\underbrace{\Bigl(Q_ky_k-\sum_j d_jx_{kj}\Bigr)}_{\ge0}\\
&+\sum_{k,j}s_{kj}\underbrace{(y_k-x_{kj})}_{\ge0}-\sum_k t_k\underbrace{y_k}_{\le1}\\
\ge{}& \Bigl(\sum_j v_j-\sum_k t_k\Bigr)+\sum_k r_ky_k+\sum_{k,j}\bar c_{kj}x_{kj}
\ \ge\ L+r_iy_i ,
\end{align*}
where the last step uses strong duality \eqref{eq:dual} and $r_k,\bar c_{kj}\ge0$. Every integer
solution is LP-feasible, so $y_i=1$ gives cost at least $L+r_i$.
\end{proof}

\noindent\emph{Certificate.} If stage $s$ solves $\mathrm{P}(K_s)$ to optimality with value $U$
and
\begin{equation}\label{eq:cert}
L+r_i\ \ge\ U\qquad\text{for all } i\in I\setminus K_s,
\end{equation}
then $U=z^\ast$: a better solution would have to open some pruned facility, which
Proposition~\ref{prop:rc} excludes. The expansion then stops with a \emph{proven} global optimum
although only a fraction of the facilities entered the solver. \emph{Guided expansion.} When
\eqref{eq:cert} fails, the set $\{i\notin K_{s}: L+r_i<U\}$ of facilities that the certificate
cannot exclude is added to the next stage, provided the stage does not exceed $1.5\,\rho_{s+1}m$
facilities. The LP is solved once (HiGHS) and its duals are reused by every later component.

With $p=p^{\mathrm{GNN}}$ this is the learned expansion; with $p=p^{\mathrm{LP}}$
\eqref{eq:lpscore} nothing is learned and the method becomes a variant of kernel search with
nested kernels, which is the natural control for the value of the GNN.

\subsection{Kernel search}\label{sec:kernel}
As the classical counterpart we implement kernel search \citep{guastaroba2014} over facilities.
The initial kernel is $\Lambda_0=\mathrm{rep}$ of $\{i: y^{\mathrm{LP}}_i>0\}$. The remaining
facilities, sorted by increasing reduced cost $r_i$, are cut into buckets
$B_1,\dots,B_q$ of size $|\Lambda_0|$. The method solves $\mathrm{P}(\Lambda_0)$ and then, for
$b=1,\dots,q$, the problem $\mathrm{P}(\Lambda_{b-1}\cup B_b)$ warm-started with the incumbent
$a$, and updates
\begin{equation}
\Lambda_b=\Lambda_{b-1}\cup\bigl(B_b\cap a(J)\bigr),
\end{equation}
i.e., a bucket facility joins the kernel only if the new incumbent uses it. The kernel problem
receives 25\% of the remaining time and each bucket an equal share of what is left. Unlike the
original, we do not add the constraint $\sum_{i\in B_b}y_i\ge1$ nor an objective cutoff; the
warm start plays the role of the cutoff. The difference from the adaptive expansion is
structural: kernel search visits \emph{disjoint} buckets around a small kernel and never solves
the full problem, whereas the expansion solves \emph{nested} sets and ends with $K=I$.

\subsection{Cluster-oriented large neighborhood search}\label{sec:clns}
\paragraph{Clusters.} Customers are clustered once by $k$-means on the log-scaled unit-cost
profile $\xi_j=\bigl(\log(1+u_{ij}/\tilde u)\bigr)_{i\in I}\in\mathbb{R}^m$, with $\tilde u$ the
median of all $u_{ij}$ and $k=\mathrm{round}(n/15)$:
\begin{equation}
\min_{G_1,\dots,G_k}\ \sum_{g=1}^{k}\sum_{j\in G_g}\lVert \xi_j-\mu_g\rVert_2^2,\qquad
\mu_g=\frac{1}{|G_g|}\sum_{j\in G_g}\xi_j .
\end{equation}
Two customers are close when they see the facilities at similar prices, which is the relevant
notion of proximity for exchanging them between facilities and does not require coordinates
(the Holmberg instances have none). The distance between clusters $g$ and $h$ is the
$\ell_1$ distance between their mean unit-cost profiles.

\paragraph{Subproblem.} Given a feasible assignment $a$ and a set $F\subseteq J$ of \emph{free}
customers, let $\bar F=J\setminus F$,
\begin{equation}\label{eq:resid}
\ell^{\bar F}_i=\sum_{j\in\bar F:\ a(j)=i}d_j,\qquad
\tilde Q_i=Q_i-\ell^{\bar F}_i,\qquad
I^{\mathrm{paid}}=a(\bar F),
\end{equation}
be the load of the fixed customers, the residual capacity, and the facilities whose fixed cost is
already paid by a fixed customer. Each free customer $j$ may use the set $C_j$ of its ten cheapest
facilities plus its current one $a(j)$. The subproblem $\mathrm{S}(a,F)$ is
\begin{align}
\min\ & \sum_{i\notin I^{\mathrm{paid}}} f_i y_i+\sum_{j\in F}\sum_{i\in C_j} c_{ij}x_{ij}
\label{eq:subobj}\\
\text{s.t.}\ & \sum_{i\in C_j}x_{ij}=1 && j\in F, \nonumber\\
& \sum_{j\in F:\ i\in C_j} d_jx_{ij}\le \tilde Q_i\,y_i && i\notin I^{\mathrm{paid}}, \nonumber\\
& \sum_{j\in F:\ i\in C_j} d_jx_{ij}\le \tilde Q_i && i\in I^{\mathrm{paid}}, \nonumber\\
& x_{ij}\le y_i && i\notin I^{\mathrm{paid}},\ j\in F,\ i\in C_j, \nonumber\\
& x_{ij},y_i\in\{0,1\}. \nonumber
\end{align}
SCIP solves it starting from $a|_F$ (always feasible, since $a(j)\in C_j$) with a short time limit
$\theta$: $\theta=0.5$~s, multiplied by $1+0.5\min\{\phi,4\}$ after $\phi$ consecutive failures.
The move is accepted only if it strictly improves the full objective \eqref{eq:za}.

\begin{proposition}[Coordination]\label{prop:feasible}
Let $a$ be feasible and $a'$ be obtained from $a$ by replacing $a|_F$ with a feasible solution of
$\mathrm{S}(a,F)$ of value $\zeta'$. Then $a'$ is feasible for
\eqref{eq:obj}--\eqref{eq:int} and
\begin{equation}\label{eq:delta}
z(a')=\zeta'+\underbrace{\sum_{i\in I^{\mathrm{paid}}}f_i+\sum_{j\in\bar F}c_{a(j)j}}_{\text{constant in }F},
\qquad\text{so}\qquad z(a)-z(a')=\zeta-\zeta',
\end{equation}
where $\zeta$ is the value of $a|_F$ in $\mathrm{S}(a,F)$.
\end{proposition}
\begin{proof}
Single sourcing holds because each free customer receives exactly one facility and fixed
customers keep theirs. For capacity, $\ell_i(a')=\ell^{\bar F}_i+\sum_{j\in F}d_jx_{ij}\le
\ell^{\bar F}_i+\tilde Q_i=Q_i$. For the cost, a facility in $I^{\mathrm{paid}}$ is open in both
$a$ and $a'$ and contributes $f_i$ once to the constant; a facility outside
$I^{\mathrm{paid}}$ is open in $a'$ iff some free customer uses it, i.e., iff $y_i=1$ in an
optimal completion of the subproblem, so its fixed cost is counted exactly once, in $\zeta'$.
\end{proof}
\noindent Proposition~\ref{prop:feasible} is what separates the CLNS from the naive
decomposition of Section~\ref{sec:coupling}: residual capacities rule out the overload
\eqref{eq:overload} and the set $I^{\mathrm{paid}}$ rules out the double counting
\eqref{eq:double}. It also shows that the improvement of the full objective equals the
improvement of the subproblem, so no global evaluation is needed to decide acceptance.

\paragraph{Candidate subproblems.} At each iteration a set $\mathcal{F}(a)$ of candidate free
sets is built from the incumbent:
\begin{itemize}
\item \emph{interior}: $F=G_g$ for each cluster $g$;
\item \emph{boundary}: $F=G_g\cup G_h$ for each cluster $g$ and its nearest cluster $h$, which
allows transfers between clusters;
\item \emph{release}: for the six open facilities with the largest fixed cost per unit of load
$f_i/\ell_i(a)$, the customers of $i$ and of the two open facilities with the closest unit-cost
profiles (subsampled to the customers of $i$ plus 30 others when the set exceeds 45
customers). This is the move that can \emph{close} a facility whose customers are spread over
several clusters, which interior and boundary moves cannot;
\item \emph{regret}: for the four customers with the largest LP regret
\begin{equation}\label{eq:regret}
\varrho_j(a)=\bar c_{a(j)j}-\min_{i\in I}\bar c_{ij}\ \ge 0,
\end{equation}
the 15 customers closest to $j$ in unit-cost profile. $\varrho_j$ is zero when $j$ is assigned as
in the LP and grows with the price the relaxation attaches to the current choice;
\item \emph{opening} (only when a GNN score is available): for the four closed facilities with the
highest $p_i$, their 15 cheapest customers.
\end{itemize}

\subsection{Choosing the subproblem}\label{sec:selectors}
A \emph{selector} maps $\mathcal{F}(a)$ and the search history to one candidate.
\emph{Rotation} cycles through the types and draws a candidate uniformly within the type.
\emph{ALNS} \citep{ropke2006} draws a type $\tau$ with probability
$\pi_\tau=\omega_\tau/\sum_{\tau'}\omega_{\tau'}$ and updates the weights every ten iterations,
\begin{equation}
\omega_\tau\leftarrow\max\bigl\{0.05,\ (1-\lambda)\,\omega_\tau+\lambda\,\bar s_\tau\bigr\},
\qquad \lambda=0.2,
\end{equation}
where $\bar s_\tau$ is the fraction of improving moves of type $\tau$ in the segment.
\emph{Dual} picks the candidate with the largest mean regret
$\frac{1}{|F|}\sum_{j\in F}\varrho_j(a)$, halved for each previous failure of that candidate.
\emph{Learned} follows learning-to-delegate \citep{li2021delegate}: a gradient-boosted regressor
$\hat g$ predicts the improvement per second of a candidate from 16 features of $(a,F)$ and the
selector takes $\arg\max_{F\in\mathcal{F}(a)}\hat g(a,F)$ (Section~\ref{sec:training}).
\emph{GNN-guided} needs no additional training and picks the candidate whose current
configuration most disagrees with the facility scores. With
$\hat p_i=\sigma(p^{\mathrm{GNN}}_i)$,
\begin{equation}\label{eq:delta-gnn}
\delta(F)=\frac{\sum_{i\in O(F)}(1-\hat p_i)+\sum_{i\in C(F)}\hat p_i}{|O(F)|+|C(F)|},
\qquad
F^\ast=\arg\max_{F\in\mathcal{F}(a)}\ \delta(F)\,2^{-\varphi(F)},
\end{equation}
where $O(F)=a(F)$ are the open facilities serving $F$, $C(F)$ are the closed facilities among the
three cheapest of each customer in $F$, and $\varphi(F)$ counts the failures of $F$ since the
last improvement. The first sum is large when $F$ is served by facilities the model considers
unlikely; the second when likely facilities nearby are closed.

\subsection{Subproblem memory}\label{sec:memory}
An LNS that accepts only improvements can select the same subproblem repeatedly from an
unchanged neighborhood. By \eqref{eq:resid}--\eqref{eq:subobj}, $\mathrm{S}(a,F)$ depends on $a$
only through $a|_F$ (which fixes the sets $C_j$ and the starting solution) and through the fixed
loads $(\ell^{\bar F}_i)_{i\in I}$ (which fix $\tilde Q$ and $I^{\mathrm{paid}}=\{i:\ell^{\bar F}_i>0\}$).
We therefore define the key
\begin{equation}\label{eq:key}
\kappa(a,F)=H\bigl(F,\ a|_F,\ (\ell^{\bar F}_i)_{i\in I}\bigr),
\end{equation}
with $H$ a 96-bit hash, and keep a table $M:\kappa\mapsto\theta_{\max}$ of the largest time limit
with which the subproblem failed to improve.

\begin{proposition}\label{prop:key}
If $\bigl(F,a|_F,\ell^{\bar F}\bigr)=\bigl(F,b|_F,\ell'^{\bar F}\bigr)$ for two feasible assignments
$a$ and $b$, then $\mathrm{S}(a,F)$ and $\mathrm{S}(b,F)$ are the same mixed-integer program with
the same starting solution.
\end{proposition}
\noindent The proof is by inspection of \eqref{eq:resid}--\eqref{eq:subobj}. Note that the key is
\emph{coarser} than the full assignment: two assignments that differ in how fixed customers are
distributed, but produce the same fixed loads, share the key. With memory enabled, candidates
with $M[\kappa(a,F)]\ge\theta$ are removed from $\mathcal{F}(a)$ before selection (unless all are
removed); a failure with a larger limit $\theta$ is still allowed, since the solver may improve
with more time. This is the transposition table of game-tree search applied to neighborhoods:
it does not change which solutions are reachable, only avoids paying twice for the same negative
answer. With memory disabled we still compute the key of the selected candidate to
\emph{measure} the repetition rate, reported in Section~\ref{sec:res-controls}.

\subsection{The hybrid}\label{sec:hybrid}
Algorithm~\ref{alg:hybrid} runs the adaptive expansion for a share $\alpha$ of the budget
($\alpha=0.5$, fixed before the experiments) and the CLNS for the rest, from the expansion's
solution, reusing the LP duals and, for the GNN-guided selector, the facility scores.

\begin{algorithm}[t]
\caption{Hybrid: adaptive expansion followed by CLNS}\label{alg:hybrid}
\begin{algorithmic}[1]
\Require instance, budget $T$, share $\alpha$, score type (GNN or LP), selector, memory flag
\State $p \gets$ facility scores; solve strong LP $\to (L, r, \bar c)$
\State $a \gets$ \textsc{AdaptiveExpansion}$(p, L, r)$ with budget $\alpha T$
\Comment{Section~\ref{sec:expansion}}
\State clusters $\gets$ $k$-means on unit-cost profiles; $M\gets\emptyset$; $\phi\gets0$
\While{time remains}
  \State $\theta\gets0.5\,(1+0.5\min\{\phi,4\})$
  \State $\mathcal{F} \gets$ candidate subproblems of $a$ (five types)
  \If{memory} $\mathcal{F}\gets\{F\in\mathcal{F}: M[\kappa(a,F)]<\theta\}$ (if non-empty) \EndIf
  \State $F \gets$ \textsc{Select}$(\mathcal{F})$ \Comment{Section~\ref{sec:selectors}}
  \State $a' \gets$ solve $\mathrm{S}(a,F)$ for at most $\theta$ seconds
  \If{$z(a') < z(a)$} $a \gets a'$; $\phi\gets0$
  \Else\ $M[\kappa(a,F)]\gets\max\{M[\kappa(a,F)],\theta\}$; $\phi\gets\phi+1$ \EndIf
\EndWhile
\State \Return $a$ (validated against the original data)
\end{algorithmic}
\end{algorithm}

\subsection{Training the learned selector}\label{sec:training}
Labels are collected \emph{off-policy}: at states visited by a reasonable search on 60 training
instances, eight candidates drawn uniformly from $\mathcal{F}(a)$ are solved from the \emph{same}
state and, for each, the features, the improvement $\Delta_F=z(a)-z(a'_F)\ge0$ and the time
$\tau_F$ are recorded; the search then continues from the best of the eight. Solving several
candidates from the same state gives counterfactual labels that an on-policy trace (one candidate
per state) cannot provide. The regression target is the normalized improvement per second,
\begin{equation}
g(a,F)=\frac{\Delta_F/\tilde c}{\max\{\tau_F,\,0.05\}},
\end{equation}
with $\tilde c$ the median pair cost of the instance, so that targets are comparable across
instances. The 16 features are: share of customers and of demand in $F$; mean cost of $F$; mean
and maximum regret \eqref{eq:regret}; number of facilities involved; capacity slack and fixed
cost per unit of load of those facilities; attempts, cumulative gain and age of the candidate;
the capacity ratio of the instance; and the candidate type. This yielded 23{,}840 training and
6{,}256 validation labels (11.7\% with $\Delta_F>0$), split by instance. The regressor reached
AUC 0.76 for predicting $\Delta_F>0$ on held-out instances; the most important features were
capacity slack and fixed cost per unit of load.
